\chapter{Cosmological Projection: RSVP and the Persistence of Structure}
\label{ch:cosmological}

\epigraph{Reality is modeled as an ongoing process of constrained transformation
rather than a fixed geometric object in which all events coexist.}{}

\section{Against the Block Universe}

Static block-universe interpretations of general relativity 	tep{barbour1999} treat the universe as
a completed four-dimensional object in which all events coexist simultaneously in
the tenseless present of spacetime geometry. The \RSVP framework provides a dynamically
consistent alternative grounded in irreversibility, entropy gradients, and locally
constrained field relaxation.

On the \RSVP interpretation, the universe is modeled less as a completed
four-dimensional object and more as an ongoing process of constrained transformation
in which present configurations continuously reorganize into future configurations
through local interactions, thermodynamic pressures, and evolving accessibility
conditions. The entropy field $S(\mathbf{x}, t)$ is not a static property of
spacetime but a dynamical quantity whose gradients drive the evolution of the
other fields.

\section{The RSVP Lagrangian and Admissibility Geometry}

The full Lagrangian structure is developed in Appendix~\ref{app:rsvp}. The seven
governing axioms specify the coupling between the scalar density field $\Phi$, the
vector flow field $\mathbf{v}$, and the entropy field $S$. Admissible cosmological
trajectories are those satisfying the derived compatibility conditions across all
three fields simultaneously.

\subsection*{Redshift: Expansion vs. Entropic Relaxation}

The standard $\Lambda$CDM interpretation of cosmological redshift attributes the
wavelength shift of photons to metric expansion: as the universe expands, the
physical wavelength of a photon increases in proportion to the scale factor $a(t)$,
giving $z = a(t_{\rm obs})/a(t_{\rm emit}) - 1$.

The \RSVP interpretation attributes redshift to entropic differential: photons
traversing regions of entropy gradient transfer energy to the $S$ field rather
than to metric stretching. The predicted redshift for a photon traveling through
an entropy gradient $\nabla S$ over path length $L$ is approximately:
\[
  z_{\rm RSVP} \approx \int_0^L \kappa\,|\nabla S(\ell)|\,d\ell,
\]
where $\kappa$ is a coupling constant relating entropy gradient to energy loss.

The two interpretations produce observationally similar predictions at low redshift
($z \lesssim 1$) because, in the \RSVP framework, the entropy field evolves at
rates that mimic scale-factor growth over cosmic time. They diverge at high
redshift and at the level of structure formation: $\Lambda$CDM predicts that the
power spectrum of density fluctuations is seeded by inflation and modulated by
dark matter, while \RSVP predicts that structure forms through entropy gradient
focusing and scalar field condensation without requiring either an inflationary
epoch or exotic dark matter. The \RSVP framework therefore makes empirically
distinct predictions about the matter power spectrum at $k \gtrsim 0.1\,h/$Mpc
and about the Silk damping scale, which constitute its primary falsification
criteria. Cosmological redshift, on this interpretation, arises
from entropic differential rather than metric expansion: photons traversing regions
of entropy gradient lose energy to the gradient field rather than to the stretching
of space. This generates observationally similar predictions to expansion cosmology
in the low-redshift regime while making distinct predictions at the structure
formation scale.

\section{Persistence, Residue, and Structural Memory}

Physical systems do not merely evolve forward through admissible trajectories.
They accumulate structural residue: traces of prior constraint configurations that
continue influencing future evolution. Fossil manifolds---large-scale structures
whose current configuration encodes the constraint history of their formation---are
cosmological counterparts to the provenance concept introduced in
Chapter~\ref{ch:vocabulary}.

\begin{definition}{Structural Residue}
The \emph{structural residue} of a dynamical process over interval $[t_0, t]$ is
the component of the current state $\Omega_t$ that is causally attributable to
constraint configurations that no longer operate at $t$ but that produced persistent
modifications to the state space topology during $[t_0, t]$. Structural residue is
the physical instantiation of provenance: the history of the system is encoded in
its present configuration in a form that, in principle, supports reconstruction of
prior states.
\end{definition}
